% Unit 09 teaching — Intermediate Trig Integrals \& Trig-Substitution
\input{technique-note.tex}
\title{Integration Bee --- Unit 09:\\ Intermediate Trig Integrals \& Trig-Substitution}
\author{UK Uni Integration Bee prep (Jerome)\\ Official: integration.soc.srcf.net\\ Yang Topics Guide 2nd Ed}
\date{}
\begin{document}
\maketitle
\section*{Status}
Concept teaching only. Tags: \texttt{intermediate-trig}, \texttt{intermediate-trig-sub}, \texttt{yang:1.15}, \texttt{yang:1.16}.
\section{Official anchors}
\begin{itemize}
  \item Guide §1.8 Trigonometric Identities (Contents p.18).
  \item Yang §§1.15, 1.16.
\end{itemize}
\section{Learning goals}
\begin{enumerate}
  \item Use power-reduction / Pythagorean identities on $\sin^2,\tan^2$.
  \item Integrate $\sqrt{x^2+1}$ and $1/\sqrt{x^2+1}$ (hyperbolic/trig-sub).
\end{enumerate}
\section{Pedagogy}
Worked $\to$ faded $\to$ interleaved. See \texttt{pedagogy/README.md}.

\section{Worked power-reduce --- \texttt{U09-D01} (\texttt{yang:1.15})}
\[
\int_0^{\pi}\sin^2 x\,dx = \int_0^{\pi}\frac{1-\cos 2x}{2}\,dx=\frac\pi2.
\]

\section{Worked hyp/trig-sub --- \texttt{U09-D05} (\texttt{yang:1.16})}
\[
\int_0^1\frac{1}{\sqrt{x^2+1}}\,dx=\operatorname{arsinh}1=\ln(1+\sqrt2).
\]

\section{Early half-angle substitution}
The syllabus puts the substitution $t=\tan(x/2)$ in section 1. The Techniques Guide teaches it as T-sub in \S1.6. Yang files the same move in a later chapter. You do not wait for that chapter when the integrand is rational in $\sin x$ and $\cos x$.
\[
\sin x=\frac{2t}{1+t^2},\qquad
\cos x=\frac{1-t^2}{1+t^2},\qquad
dx=\frac{2\,dt}{1+t^2}.
\]
The fourth replacement is $\tan x=2t/(1-t^2)$. It is the quotient of the first two. Use it when a tangent is already in the integrand. $t=\tan(x/2)$ is continuous on $(-\pi,\pi)$. If the interval contains an odd multiple of $\pi$ in its interior, split there before you integrate. If the endpoint is $x=\pi$, the integral in $t$ is improper. Unit 26 writes that endpoint.

\section{Worked early pass}
This integral is not one of \texttt{U09-D01} through \texttt{U09-D08}. On $[0,\pi/3]$,
\[
\int_0^{\pi/3}\frac{dx}{1+\cos x}.
\]
The limits become $t=0$ and $t=\tan(\pi/6)=\sqrt{3}/3$. Then
\[
1+\cos x=\frac{2}{1+t^2},
\]
so $1/(1+\cos x)=(1+t^2)/2$. Multiply by $dx=2\,dt/(1+t^2)$ and the integral is $\int_0^{\sqrt{3}/3} dt=\sqrt{3}/3$.
The antiderivative is $\tan(x/2)$. Check the endpoint: $\tan(\pi/6)=\sqrt{3}/3$.

\section{Faded early pass --- \texttt{U26-D01}}
\[
\int_0^{\pi/2}\frac{dx}{1+\sin x}.
\]
Write $t=\tan(x/2)$ and fill the limits before you expand. They are $0$ and $1$, not $0$ and $\infty$. The identity $1+\sin x=(1+t)^2/(1+t^2)$ collapses the integral to $\int_0^1 2/(1+t)^2\,dt$. The stored value is $1$. The drill id lives in unit 26. That file is the note a miss of this substitution names.

A power of sine alone is still the reduction in unit 20. A rational function of $\tan x$ that already contains $\sec^2 x$ is the ordinary substitution $u=\tan x$, which is \texttt{U09-D04} and \texttt{U09-D08}, not this half-angle. The full pattern $\int_0^\pi dx/(a+b\cos x)$ is unit 26.

\section{Practice}
\texttt{U09-D01}\ldots\texttt{U09-D08} in \texttt{drills.tex}; \texttt{verify/run.sh}.
\end{document}
